% Q05_solutions.tex — REFERENCE SOLUTIONS to paper/Q05_problems.tex.
% Author-only and read-denied in study mode (Read(./paper/solutions/**)).
% Built only by `make -C paper paper-solutions` (scripts/ci.sh, or the
% author's linked worktree); `make paper` never touches this directory.
%
% Every verifier result quoted below was captured on the course pins (Dafny
% 4.11.0, Verus 0.2026.09.27) through scripts/vcheck, one solver job at a time,
% by the DV3k kit lane for Q05 on 2026-10-06 (the problems' code extracted
% verbatim from the problems file into instructor/scratch/probe.{dfy,rs} of a
% scratch course root; probe sources archived under
% ~/.cache/dafny-verus-learn-0/author-probes/dv3k/q05-kit/). B4's fields are
% graded against the learner's own cheatsheets.md; this file records which
% lines of the pinned range each field draws on and never copies the card.
\documentclass[11pt]{article}
\input{preamble}

\title{Q05 --- Frames, ownership, modes, wp$\to$VC \\ \large Reference solutions (author-only)}
\author{dafny-verus-learn-0}
\date{}

\begin{document}
\maketitle

\section*{How these are graded}
B1 and B2 are hand calculations, dimension \texttt{recomputation}: the grader
re-derives each step with the rubric's Probe A and Probe B and diffs the
learner's trace step by step. B3 is informal-proof-first (the five proof
dimensions); B4 is a paper card (\texttt{card-fidelity}, per field). The
solutions below give the answers; the rubric never states them.

% ---------------------------------------------------------------------------
\section*{B1 (hand-calculation)}

Write $L$ for \dfy{a.Length}, $s$ for the contents at the start, $c$ for the
contents after the first call and $c'$ after the second.
$R \equiv t = s[0] + s[1] \wedge c'[0] = 0 \wedge c'[1] = 0$. By the sequence
rule, $\WP{\mathit{body}}{R} = \WP{x \assign \mathit{Drain}(a,0)}{\WP{y \assign
\mathit{Drain}(a,1)}{\WP{t \assign x + y}{R}}}$.

\textbf{(a)}
\begin{enumerate}
\item Assignment: $R_1 \equiv \subst{R}{t}{x + y} \equiv x + y = s[0] + s[1]
  \wedge c[0] = 0 \wedge c[1] = 0$, over the contents $c$ before the
  assignment.
\item Call \dfy{y := Drain(a, 1)}, from contents $c$:
  \[ \begin{aligned} R_2 = {} & 1 < L \wedge \forall c', y' .\; \bigl(y' = c[1] \wedge c'[1] = 0
     \wedge \forall j .\; 0 \leq j < L \wedge j \neq 1 \impl c'[j] = c[j]\bigr) \\
     & \impl \bigl(x + y' = s[0] + s[1] \wedge c'[0] = 0 \wedge c'[1] = 0\bigr). \end{aligned} \]
  Simplify. The first hypothesis fixes $y' = c[1]$; the second gives the third
  conclusion; the frame hypothesis at $j = 0$ (in range: $0 < 1 < L$) fixes
  $c'[0] = c[0]$. The hypotheses are satisfiable ($c'$ = $c$ with limb 1 set to
  $0$), so the $\forall$ is equivalent to its conclusion at that outcome:
  \[ R_2 \equiv 1 < L \wedge x + c[1] = s[0] + s[1] \wedge c[0] = 0. \]
\item Call \dfy{x := Drain(a, 0)}, from the start contents $s$:
  \[ \begin{aligned} R_3 = {} & 0 < L \wedge \forall c, x' .\; \bigl(x' = s[0] \wedge c[0] = 0
     \wedge \forall j .\; 0 \leq j < L \wedge j \neq 0 \impl c[j] = s[j]\bigr) \\
     & \impl \bigl(1 < L \wedge x' + c[1] = s[0] + s[1] \wedge c[0] = 0\bigr). \end{aligned} \]
  Simplify. $1 < L$ does not mention $c$ or $x'$ and the hypotheses are
  satisfiable, so it leaves the $\forall$. Under $1 < L$ the frame hypothesis at
  $j = 1$ fixes $c[1] = s[1]$, and $x' = s[0]$, so $x' + c[1] = s[0] + s[1]$;
  $c[0] = 0$ is a hypothesis. Hence
  \[ \WP{\mathit{body}}{R} \equiv 0 < L \wedge 1 < L \;\equiv\; L \geq 2, \]
  with no constraint on $s$.
\end{enumerate}
The two call steps are the chapter's four steps each: prove the precondition,
forget the frame, assume the postcondition, continue. Simplifying only at the
end gives the same result if each step is an equivalence.

\textbf{(b)} The verification condition is $L = 2 \impl L \geq 2$: valid.
Without the frame postcondition, step 2 loses the hypothesis that fixed
$c'[0]$, and step 3 the one that fixed $c[1]$:
\[ \begin{aligned} R_2' = {} & 1 < L \wedge \forall c', y' .\; (y' = c[1] \wedge c'[1] = 0) \\
   & \impl (x + y' = s[0] + s[1] \wedge c'[0] = 0 \wedge c'[1] = 0), \end{aligned} \]
which is false (an outcome with $c'[0] = 7$ is allowed), and so is the wp: the
verification condition fails for every start state. Two conjuncts of $R$ can
no longer be proved: $t = s[0] + s[1]$ and $c'[0] = 0$; the third,
$c'[1] = 0$, is the second call's own postcondition and still holds.
Instance: $s = [3, 4]$; the first call returns $x = 3$ and leaves $c = [0, 9]$
(allowed: $x = s[0]$, $c[0] = 0$); the second returns $y = 9$ and leaves
$c' = [7, 0]$ (allowed: $y = c[1]$, $c'[1] = 0$). Then $t = 12 \neq 7$ and
$c'[0] = 7 \neq 0$. ($c' = [0, 0]$ instead breaks the sum alone.) No run ends
that way: \dfy{Drain}'s body writes only \dfy{a[i]}, and every run from
$[3, 4]$ returns $7$ and leaves $[0, 0]$. The failure is about the contract,
not the program: a caller knows a call only through its contract, and the
deleted \dfy{ensures} was all that told \dfy{Collect} that limb 1 survives the
first call and limb 0 the second. On the pins, the Dafny program with that
line deleted gives two errors, both \texttt{a postcondition could not be
proved on this return path} at \dfy{Collect}'s body, related to the first and
the second conjunct of its \dfy{ensures} (columns 13 and 46); the program as
given verifies (\texttt{4 verified}).

\textbf{(c)} (i) The \vrs{&mut} parameter: \vrs{v: &mut Vec<u64>} lends
\vrs{*v}, all of it, length included, and nothing else the caller holds; the
\vrs{forall} clause (with the length clause) is what the Dafny frame
postcondition is, the part of \vrs{*v} that stayed. (ii) The overflow
obligation of \vrs{x + y} on \vrs{u64}: $x + y \leq 2^{64} - 1$. It rests on
\vrs{collect}'s bounds \vrs{old(v)@[0] < 1000} and \vrs{old(v)@[1] < 1000},
on \vrs{drain}'s \vrs{x == old(v)@[i as int]} for each call, and on the
\vrs{forall} clause of the first call, which carries limb 1 to the second call
unchanged. Dafny's \dfy{nat} addition has no such obligation. (iii) With the
\vrs{forall} clause deleted, two obligations fail: the overflow check
(\texttt{possible arithmetic underflow/overflow} at \vrs{x + y}) and the
postcondition. Verus reports one \texttt{postcondition not satisfied} per
return path, at the first failing clause (\vrs{t == old(v)@[0] + old(v)@[1]});
\vrs{final(v)@[0] == 0} fails too (a probe without the first clause reports
it), and \vrs{final(v)@[1] == 0} holds. As given, \vrs{collect} verifies
(\texttt{2 verified}).

\textbf{Probe A on the pins.} With $L \geq 2$ spliced into \texttt{LearnerWp},
Probe A's \texttt{Sufficient} verifies (\texttt{4 verified}); in Probe A$''$,
\texttt{Included} on $[3, 4]$ passes and \texttt{Excluded} on a one-limb array
fails at \texttt{Drain's requires fails}: neither state is a counterexample.
With a formula too strong ($L = 2 \wedge s[0] > 0$, guarded by $|s| = 2$),
\texttt{Excluded} on $[0, 5]$ passes: a refutation of necessity. With
\texttt{true}, \texttt{Sufficient} fails at both calls' preconditions and at
the postcondition's indexing. Probe A$'$ on step 3, with the direction from
the quantified formula naming the outcome $s[0 \mapsto 0]$ (\dfy{s[0 := 0]}),
verifies in both directions (\texttt{4 verified}); with $L \geq 1$ as the
result, \texttt{StepUp} fails. Probe A$'''$ with the instance of (b) passes
(\texttt{PASSED WeakenedRun}) for each of the two conjuncts, and fails at
\texttt{that conjunct of R holds} for $c'[1] = 0$.

\textbf{Typical flaws, with their failing instances.}
\begin{itemize}
\item \emph{The body read instead of the contract} in (b): ``the VC still holds,
  since \dfy{Drain} only writes \dfy{a[i]}''. Failing instance: the outcome
  above ($s = [3, 4]$, $c = [0, 9]$, $c' = [7, 0]$), allowed by the weakened
  contract, on which $R$ fails; Probe A$'''$ passes on it.
\item \emph{A precondition dropped} in step 3 (the result $0 < L$, i.e.\
  $L \geq 1$): a one-limb array satisfies the claimed wp, and the second
  call's precondition $1 < L$ fails on it: Probe A$''$'s \texttt{Included}
  on it fails at \texttt{Drain's requires fails} (and Probe A's
  \texttt{Sufficient} reports the second call).
\item \emph{A result over an intermediate state} ($L \geq 2 \wedge c[1] = s[1]$,
  which names the contents after the first call): the wp is a formula over
  the start state, so the simplification of step 3 is unfinished; a gap (3),
  unless the formula is false at an exhibited start state.
\end{itemize}

% ---------------------------------------------------------------------------
\section*{B2 (hand-calculation)}

$I \equiv 0 \leq i \leq L \wedge \forall j .\; 0 \leq j < i \impl c[j] = 0$,
$b \equiv i < L$, $S \equiv x \assign \mathit{Drain}(a, i) \seqc i \assign i + 1$,
measure $L - i$.

\textbf{(a)}
\begin{itemize}
\item \emph{Entry:} $\WP{i \assign 0}{I} = 0 \leq 0 \leq L \wedge \forall j .\;
  0 \leq j < 0 \impl c[j] = 0$. Valid: $L \geq 0$ for an array, and the range
  is empty.
\item \emph{Maintenance:} $\WP{i \assign i + 1}{I} = 0 \leq i + 1 \leq L \wedge
  \forall j .\; 0 \leq j < i + 1 \impl c[j] = 0$; then the call rule:
  \[ \begin{aligned} \WP{S}{I} = {} & i < L \wedge \forall c', x' .\; \bigl(x' = c[i] \wedge c'[i] = 0
     \wedge \forall j .\; 0 \leq j < L \wedge j \neq i \impl c'[j] = c[j]\bigr) \\
     & \impl \bigl(0 \leq i + 1 \leq L \wedge \forall j .\; 0 \leq j < i + 1 \impl
     c'[j] = 0\bigr). \end{aligned} \]
  $I \wedge b \impl \WP{S}{I}$ is valid: $i < L$ is the guard and gives
  $i + 1 \leq L$; for $j < i$, the frame gives $c'[j] = c[j]$ and $I$ gives
  $c[j] = 0$; for $j = i$, $c'[i] = 0$ by the postcondition.
\item \emph{Use at exit:} $I \wedge \neg b \impl \forall j .\; 0 \leq j < L
  \impl c[j] = 0$. Valid: $i \leq L$ and $i \geq L$ give $i = L$.
\item \emph{Termination:} bounded, $I \wedge b \impl L - i \geq 0$ (indeed
  $> 0$); decreasing, $I \wedge b \impl \WP{S}{L - i < V_0}$ with $V_0 = L - i$
  at the start of the iteration, that is $i < L \wedge \forall c', x' .\;
  (\dots) \impl L - (i + 1) < L - i$. Valid: the call cannot change $L$, and
  $i$ goes up by one.
\end{itemize}

\textbf{(b)} Without the frame postcondition, entry, use at exit and
termination are unchanged (none uses the frame) and stay valid.
\emph{Maintenance fails}: the call outcome is constrained only by
$x' = c[i] \wedge c'[i] = 0$, so $c'[j] = 0$ for $j < i$ no longer follows.
Instance: $L = 2$, $i = 1$, $c = [0, 5]$ (it satisfies $I \wedge b$), and the
outcome $x' = 5$, $c' = [9, 0]$; the conclusion needs $c'[0] = 0$. No run
reaches it: \dfy{Drain}'s body writes only \dfy{a[i]}. No invariant repairs
it: the postcondition needs every limb $0$ after the last call, and the
weakened contract lets that call (at $i = L - 1$) leave any value in every
other limb, so nothing an invariant says before it survives it (for
$L \geq 2$). The failure is about the callee's contract, which is too weak for
its callers; in Q03's B2 the program and the code were fine and the
\emph{invariant} was too weak (an unreachable state satisfied it). Both are
verification conditions that fail with a correct program; what differs is
which annotation is too weak, and so where the repair goes: here the
\dfy{ensures} of \dfy{Drain}, there the invariant. On the pins: as given,
\texttt{4 verified}; with the line deleted, one error, \texttt{this invariant
could not be proved to be maintained by the loop}, at the second invariant.

\textbf{(c)} Two conditions fail: maintenance of \vrs{i <= v.len()} (after
the call the length is unknown) and termination (the measure
\vrs{v.len() - i} reads the length the call may now change). On the pins:
\texttt{invariant not satisfied at end of loop body} at \vrs{i <= v.len()}
and \texttt{decreases not satisfied at end of loop}; as given, \vrs{drain_all}
verifies (\texttt{3 verified}). The Dafny \dfy{Drain} needs no such line: an
array's length is fixed when the array is allocated, and \dfy{modifies a} lets
the callee change \dfy{a}'s elements, never \dfy{a.Length}. A Verus
\vrs{Vec}'s length is part of its value, and \vrs{&mut v} lends all of it. (With
\vrs{drain}'s \vrs{forall} clause deleted instead, the \vrs{forall} invariant
fails to be maintained, as in (b).)

\textbf{Typical flaws.} Maintenance computed without the call rule (treating
the call as an assignment to \dfy{a[i]} alone, which is reading the body);
the use condition written with $b$ instead of $\neg b$; (b) answered ``use at
exit fails'' (Probe B's \texttt{UseAtExit} verifies, untouched by the
weakened contract: a verified refutation of the claim); (b) answered ``a stronger invariant
fixes it'' (the outcome above refutes any invariant: the last call may break
limb 0); (c) answered with the \vrs{forall} invariant (Probe B$''$'s verdict
names \vrs{i <= v.len()} and the measure). On the pins, Probe B filled with
the four conditions of (a) verifies (\texttt{8 verified}); with the weakened
contract in the call's \dfy{requires}, only \texttt{Maintenance} fails, and
Probe B$'$ on the state of (b) passes (\texttt{PASSED WeakenedMaintenance}).

% ---------------------------------------------------------------------------
\section*{B3 (informal-proof-first): the reference blueprint}

\textbf{Statement.} For every sequence $s$ of naturals, the loop ends with
$c = \mathrm{Ones}(s)$ (and, in Verus, under $\mathrm{ones}(v) \leq 2^{64} - 1$,
no arithmetic overflows).

\textbf{Loop invariant and measure.}
\[ I \;\equiv\; 0 \leq i \leq |s| \;\wedge\; c + \mathrm{BitCount}(x) +
   \mathrm{Ones}(s[i..]) = \mathrm{Ones}(s), \]
\[ d \;\equiv\; (|s| - i,\; x) \text{ compared left to right.} \]
The count so far, the bits still in $x$ and the bits of the limbs not yet
loaded add up to the total.

\textbf{Entry.} $i = 0$, $x = 0$, $c = 0$: $\mathrm{BitCount}(0) = 0$ and
$s[0..] = s$, so $0 + 0 + \mathrm{Ones}(s) = \mathrm{Ones}(s)$.

\textbf{Maintenance} (the guard is $x > 0 \vee i < |s|$), one case per branch:
\begin{enumerate}
\item $x > 0$: $c' = c + x \bmod 2$, $x' = x / 2$, $i$ unchanged. One
  unfolding of BitCount ($x \neq 0$): $\mathrm{BitCount}(x) = x \bmod 2 +
  \mathrm{BitCount}(x/2)$, so $c' + \mathrm{BitCount}(x') = c +
  \mathrm{BitCount}(x)$ and the sum is unchanged. Measure: first component
  unchanged, $x/2 < x$ because $x > 0$.
\item $x = 0$, so $i < |s|$ by the guard: $x' = s[i]$, $i' = i + 1$, $c$
  unchanged. $s[i..]$ is non-empty, so one unfolding of Ones gives
  $\mathrm{Ones}(s[i..]) = \mathrm{BitCount}(s[i]) + \mathrm{Ones}(s[i..][1..])$,
  and the given assertion $s[i..][1..] = s[i+1..]$ turns the tail into
  $s[i+1..]$. With $\mathrm{BitCount}(0) = 0$:
  $c + \mathrm{BitCount}(s[i]) + \mathrm{Ones}(s[i+1..]) = c + 0 +
  \mathrm{Ones}(s[i..])$. $0 \leq i + 1 \leq |s|$. Measure: the first
  component drops by one; $x$ may grow, which a lexicographic order allows.
\end{enumerate}
Bounded below: both components are naturals ($|s| - i \geq 0$ by $I$).
Neither component alone goes down at every pass, which is why the pair is
needed.

\textbf{Use at exit.} $\neg(x > 0 \vee i < |s|)$ gives $x = 0$ and $i \geq |s|$,
so $i = |s|$ by $I$; $\mathrm{BitCount}(0) = 0$ and $\mathrm{Ones}(s[|s|..])
= \mathrm{Ones}([]) = 0$, so $c = \mathrm{Ones}(s)$.

\textbf{Overflow in H4v.} The update $c + x \bmod 2$ happens only when
$x > 0$, where $\mathrm{bit\_count}(x) = x \bmod 2 + \mathrm{bit\_count}(x/2)
\geq x \bmod 2$. So $c + x \bmod 2 \leq c + \mathrm{bit\_count}(x) \leq
\mathrm{ones}(v) \leq 2^{64} - 1$ by $I$ and the precondition, which the loop
must carry as an invariant (a loop sees only its invariants). Loading
\vrs{v[i]} into the \vrs{u64} $x$ and $i + 1 \leq$ \vrs{v.len()} cannot
overflow.

\textbf{Helper facts} (no new lemma): one unfolding of BitCount at $x > 0$ and
its value at $0$; one unfolding of Ones on a non-empty sequence and its value
on the empty one; the given $s[i..][1..] = s[i+1..]$; $s[0..] = s$ at entry
(Dafny proves it unaided; in Verus it is an extensional equality,
\vrs{v@.subrange(0, len) =~= v@}, asserted before the loop, and the load
step's tail fact is the \vrs{drop_first} of the suffix asserted with
\vrs{=~=}).

\textbf{Gaps:} none.

\textbf{Structure match (spec \S4.4 author rule).} The formal references
follow this blueprint: the Dafny twin's H4d region is the two invariants and
\dfy{decreases |s| - i, x}; the Verus twin's \vrs{count_ones} keeps the same
invariant over \vrs{v@.subrange(i, len)}, carries
\vrs{ones(v@) <= u64::MAX} as a third invariant, asserts the entry equality
and the load step's \vrs{drop_first} fact with \vrs{=~=}, and decreases
\vrs{v.len() - i, x}.

\textbf{Probe C on the pins.} The problem's code with these two invariants and
this measure spliced in verifies under \texttt{scripts/vcheck --module Q05}
(\texttt{25 verified}, the spine's tasks included; \texttt{3 verified} alone).
With the measure reversed (\dfy{decreases x, |s| - i}) or shortened to
\dfy{decreases |s| - i}: \texttt{decreases expression might not decrease}.
With the BitCount term dropped from the invariant: \texttt{this invariant
could not be proved to be maintained by the loop}. Probe C$''$'s overflow lemma,
with the facts above as its hypotheses, verifies (\texttt{25 verified} under
\texttt{--module Q05}); with $c + \mathrm{rest} \leq \mathrm{total}$ in place of
the equation it does not.

\textbf{Typical flaws.}
\begin{itemize}
\item \emph{The measure in the wrong order} ($x$ first, claimed to drop when a
  limb is loaded): induction-soundness 2. Failing instance: $s = [5]$,
  $i = 0$, $x = 0$, $c = 0$: the load pass gives $(x, |s| - i)$ from $(0, 1)$
  to $(5, 0)$, larger in the first component; Probe C$'$'s one-pass test passes
  with the decrease negated.
\item \emph{One component only} ($|s| - i$): the $x > 0$ branch leaves it
  unchanged. Instance: $s = [2]$, $i = 1$, $x = 2$.
\item \emph{The bits of $x$ left out of the invariant}
  ($c + \mathrm{Ones}(s[i..]) = \mathrm{Ones}(s)$): not maintained at the load
  step. Instance: $s = [1]$, $i = 0$, $x = 0$, $c = 0$ satisfies it; after the
  load, $0 + \mathrm{Ones}([]) = 0 \neq 1$.
\item \emph{The overflow paragraph says ``$c \leq \mathrm{ones}(v)$''} without
  the step $x \bmod 2 \leq \mathrm{bit\_count}(x)$: true, but the step is
  unjustified, a gap (3).
\end{itemize}

% ---------------------------------------------------------------------------
\section*{B4 (paper card 03.12): reference fields and their lines}

Pinned range: \texttt{cheatsheets.md} lines 1366--1386 (computed 2026-10-06),
\[ \texttt{sha256 a821f639a8a72dc663372dad601c93a6ea0fbc8cee85faf9421669c8ac521414}. \]
A field is correct when it agrees with the lines named; wording is free. Each
entry: the reference content (paraphrased), then the lines it draws on.

\begin{description}
\item[Task shape.] Input: a small C program with a property to prove, a
  predicate at a program line. Output: a proof that the property holds on
  every execution (success) or that it does not (fail), found by a calculus of
  eight rules over configurations (program, at most one assumption, a trail
  of goals) in which a language model proposes loop invariants and sub-goals
  as assertions at marked lines. \emph{Lines 1368, 1370, 1373, 1376.}
\item[Oracle.] A conventional C verifier (esbmc or UAutomizer), sound but
  incomplete (True, False or Unknown). Every step is decided by it: a
  proposal is used only after the verifier proves the implication it is
  meant for, and the calculus's theorems make success imply the property and
  fail imply its negation, whatever the model proposes. \emph{Lines 1368,
  1370, 1373, 1374.}
\item[What is trusted.] Soundness of the verifier and of the calculus: its
  two soundness theorems, and Lemma 2.1's case split, which needs a stable
  assumption; the stability check itself, which the paper does not describe.
  Not the language model, whose output only steers speed and completeness.
  \emph{Lines 1368, 1373, 1374, 1384, 1386.}
\item[Cheat surface.] Closed: a wrong or vacuous proposal cannot produce a
  verdict, because the checker decides every implication (the canonical trust
  boundary). Open: the unspecified stability check; a hard set whose
  properties are all true, so the fail rule is never exercised there, and
  which was selected because the baselines fail on it; possible contamination
  of the tiny programs. \emph{Lines 1373, 1374, 1384, 1386.}
\item[Metric.] Benchmarks solved within a wall-clock limit: Code2Inv (133, 10
  minutes) 103 with GPT-3.5 and 107 with GPT-4, against esbmc's 68 and
  Code2Inv's quoted 92 (at 1 hour); a hard SV-COMP subset (47, 15 minutes)
  14--15 and 24--25 (19--20 without repair), against 1 for UAutomizer and
  esbmc (2 and 4 at 12 hours). \emph{Lines 1379--1382.}
\item[Budget.] A budget of $k$ proposals per goal (needed for termination,
  with every proposal at an earlier line), 30 seconds per verifier call, and
  10 or 15 minutes per benchmark; no token, call-count or dollar figure (token
  limits appear only as a limitation). \emph{Lines 1375, 1376, 1381, 1382,
  1384.}
\item[Variance.] Each configuration ran 3 times; the hard-set results are
  ranges over those runs (14--15, 24--25, 19--20); no intervals or seeds are
  reported. \emph{Lines 1382, 1384.}
\item[Artifact.] The calculus with its eight rules and Algorithm 1,
  instantiated with GPT-3.5 or GPT-4 and esbmc or UAutomizer, evaluated on
  Code2Inv and a hard SV-COMP subset; the range does not say whether code or
  prompts are released. \emph{Lines 1368, 1373, 1375, 1379--1382.}
\end{description}
A field that says the range is silent where it is silent (a code release; a
token or dollar budget; intervals) is correct. A field that contradicts a line
(for example ``each configuration ran once'', or ``the model's invariants are
trusted when the verifier times out'') scores 2, quoting that line.

\end{document}
